% O014 / D80 — Methods of Algebra, Volume 2 — id-ID
% Sumber: chapter9.tex, baris 1724--1804 (inklusif)
% ID stabil unit: o014.aljabr2.chapter9.fiber-functor-isomorphisms
% Koreksi terverifikasi: omega'_1 pada baris sumber 1753 menjadi omega_1;
% omega_2(X) pada baris sumber 1767 menjadi omega_1(Y), sesuai diagram.
Terakhir, kita tinjau secara singkat hubungan antarfunktor serat. Berbeda dari
pembahasan sebelumnya, bagian ini tidak melibatkan teorema rekonstruksi.
Misalkan $\mathcal{T}$ kategori Tannakian dan $\omega_1,\omega_2$ funktor serat
di atas aljabar-$\Bbbk$ komutatif $B$. Untuk setiap aljabar-$B$ komutatif
$B'$, definisikan himpunan
\begin{equation*}
	\iHom^\otimes_B(\omega_2, \omega_1)(B') := \Hom_{\text{funktor monoidal}}\left( \omega_2 \dotimes{B} B', \omega_1 \dotimes{B} B' \right).
\end{equation*}
Ketika $B'$ bervariasi, ini memberikan funktor $B\dcate{CAlg}\to\cate{Set}$. Ingat bahwa Proposisi \ref{prop:L-bialgebra} (iii) memberikan struktur aljabar-$B$ komutatif pada $\coHom_B(\omega_1,\omega_2)$.

\begin{proposition}\label{prop:Isom-Tannakian}
	Untuk kategori Tannakian $\mathcal{T}$ dan funktor serat $\omega_1,\omega_2:\mathcal{T}\to\cated{Mod}B$, terdapat isomorfisme funktor
	\[ \iHom^\otimes_B(\omega_2, \omega_1) \simeq \Hom_{B\dcate{CAlg}}\left( \coHom_B(\omega_1, \omega_2) , \; \cdot \; \right). \]
	Komposisi homomorfisme pada ruas kanan diinduksi oleh \eqref{eqn:L-comult-prep}.
\end{proposition}
\begin{proof}
	Untuk meringankan notasi, selanjutnya tuliskan $\coHom_B(\omega_1,\omega_2)$
	sebagai $\coHom$. Menentukan suatu morfisme (dengan mengabaikan struktur
	monoidal)
	\[ \tilde{\varphi}: \omega_2 \dotimes{B} B' \to \omega_1 \dotimes{B} B' \quad (\text{keduanya bernilai dalam }\cated{Mod}B') \]
sama dengan menentukan suatu morfisme
	\[ \varphi: \omega_2 \to \omega_1 \dotimes{B} B' \quad (\text{keduanya bernilai dalam }\cated{Mod}B), \]
dan, berdasarkan sifat universal $\coHom$, ini sama dengan menentukan homomorfisme modul-$B$
	\[ \psi: \coHom \to B'. \]
Masalahnya direduksi ke pembuktian bahwa $\tilde{\varphi}$ mempertahankan struktur monoidal jika dan hanya jika $\psi$ merupakan homomorfisme aljabar-$B$. Argumennya formal, meskipun cukup panjang jika dituliskan dengan jelas.
	
	Nyatakan morfisme perkalian dengan $m:B'\dotimes{B}B'\to B'$ dan $\mu:\coHom\dotimes{B}\coHom\to\coHom$. Mudah dilihat bahwa $\tilde{\varphi}$ mempertahankan unit (atau mempertahankan $\otimes$) jika dan hanya jika diagram kiri (atau kanan) berikut selalu komutatif:
	\[\begin{tikzcd}
		\omega_2(\munit) \arrow[dd, "\sim"' sloped] \arrow[r, "{\varphi_{\munit}}"] & \omega_1(\munit) \dotimes{B} B' \arrow[d, "\sim" sloped] \\
		& B \dotimes{B} B' \arrow[d, "\sim" sloped] \\
		B \arrow[r] & B'
	\end{tikzcd} \quad
	\begin{tikzcd}
		\omega_2(X \otimes Y) \arrow[r, "{\varphi_{X \otimes Y}}"] \arrow[dd, "\sim"' sloped] & \omega_1(X \otimes Y) \dotimes{B} B' \arrow[d, "\sim" sloped] \\
		& \omega_1(X) \dotimes{B} \omega_1(Y) \dotimes{B} B' \\
		\omega_2(X) \dotimes{B} \omega_2(Y) \arrow[r, "{\varphi_X \otimes \varphi_Y}"'] & \omega_1(X) \dotimes{B} B' \dotimes{B} \omega_1(Y) \dotimes{B} B' . \arrow[u, "{\text{diinduksi oleh }m}"']
	\end{tikzcd}\]

	Mula-mula tangani diagram kiri. Baris pertamanya berfaktor sebagai $\omega_2(\munit)\to\omega_1(\munit)\dotimes{B}\coHom\xrightarrow{\identity\otimes\psi}\omega_1(\munit)\dotimes{B}B'$. Setelah semua $\omega_i(\munit)$ diidentifikasi dengan $B$, komutativitas diagram kiri ekuivalen dengan komutativitas
	$\begin{tikzcd}[column sep=small, row sep=small]
		B \arrow[r] \arrow[d] & B' \\
		\coHom \arrow[ru, "\psi"'] &
	\end{tikzcd}$,
yang juga ekuivalen dengan $\psi$ mempertahankan unit.
	
	Sekarang kita tunjukkan bahwa komutativitas diagram kanan ekuivalen dengan $\psi$ mempertahankan perkalian. Perhatikan
	\begin{equation}\label{eqn:Isom-Tannakian-aux0}
		\begin{tikzcd}
			\omega_2(X \otimes Y) \arrow[rr, bend left, "{\varphi_{X \otimes Y}}"] \arrow[r] \arrow[dd, "\sim"' sloped] & \omega_1(X \otimes Y) \dotimes{B} \coHom \arrow[r, "{\identity \otimes \psi}"] & \omega_1(X \otimes Y) \dotimes{B} B' \arrow[d, "\sim" sloped] \\
			& \omega_1(X \otimes Y) \dotimes{B} \coHom \dotimes{B} \coHom \arrow[u, "{\identity \otimes \mu}"] & \omega_1(X) \dotimes{B} \omega_1(Y) \dotimes{B} B' \\
			\omega_2(X) \dotimes{B} \omega_2(Y) \arrow[rr, bend right, "{\varphi_X \otimes \varphi_Y}"'] \arrow[r] & \omega_1(X) \dotimes{B} \coHom \dotimes{B} \omega_1(Y) \dotimes{B} \coHom \arrow[r, "{(\identity \otimes \psi)^{\otimes 2}}" inner sep=0.6em] \arrow[u, "\sim" sloped] & \omega_1(X) \dotimes{B} B' \dotimes{B} \omega_1(Y) \dotimes{B} B' , \arrow[u, "{\text{diinduksi oleh }m}"']
		\end{tikzcd}
	\end{equation}
Dua daerah melengkung di atas dan bawah selalu komutatif, sedangkan komutativitas persegi di kiri direduksi ke pencirian $\mu$ dalam \eqref{eqn:L-bialgebra-prod}.
	
	Definisikan $\eta:=\omega_1^\vee\dotimes{B}\omega_2$, yang merupakan funktor monoidal $\mathcal{T}\to\cated{Mod}B$. Bagian persegi dari diagram \eqref{eqn:Isom-Tannakian-aux0} dapat disusun ulang menjadi
	\begin{equation}\label{eqn:Isom-Tannakian-aux}
		\begin{tikzcd}
			\eta(X \otimes Y) \arrow[r] \arrow[d, "\sim"' sloped] & \coHom \arrow[r, "\psi"] & B' \\
			\eta(X) \dotimes{B} \eta(Y) \arrow[r] & \coHom \dotimes{B} \coHom \arrow[r, "{\psi \otimes \psi}"'] \arrow[u, "\mu"] & B' \dotimes{B} B' , \arrow[u, "m"']
		\end{tikzcd}
	\end{equation}
Persegi di kiri tetap komutatif, sedangkan $\psi$ mempertahankan perkalian jika dan hanya jika persegi di kanan komutatif.
	
	Jika $\tilde{\varphi}$ mempertahankan $\otimes$, maka bingkai luar besar pada \eqref{eqn:Isom-Tannakian-aux0} komutatif, sehingga bingkai luar bagian perseginya juga komutatif. Jadi bingkai luar \eqref{eqn:Isom-Tannakian-aux} komutatif. Ini menunjukkan bahwa persegi kanan komutatif setelah didahului oleh $\eta(X)\dotimes{B}\eta(Y)\to\coHom\dotimes{B}\coHom$. Karena $X$ dan $Y$ sembarang, sifat universal $\coHom$, atau konstruksi konkretnya dalam Proposisi \ref{prop:cogebra-L-universal}, cukup untuk menunjukkan bahwa persegi kanan \eqref{eqn:Isom-Tannakian-aux} komutatif.
	
	Sebaliknya, andaikan persegi kanan \eqref{eqn:Isom-Tannakian-aux} komutatif. Maka seluruh diagram \eqref{eqn:Isom-Tannakian-aux} komutatif, sehingga bingkai luar besar \eqref{eqn:Isom-Tannakian-aux0} komutatif; dengan demikian $\tilde{\varphi}$ mempertahankan $\otimes$. Pemeriksaan selesai.
\end{proof}

Perhatikan bahwa kekakuan $\mathcal{T}$ menyiratkan bahwa setiap unsur $\iHom^\otimes_B(\omega_2,\omega_1)(B')$ merupakan isomorfisme; ini merupakan penerapan Proposisi \ref{prop:monoidal-functor-rigidity}. Karena itu, hasil berikut jelas.

\begin{corollary}\label{prop:Aut-Tannakian}
	\index[sym1]{Aut-tensor@$\Aut^{\otimes}$}
	Untuk kategori Tannakian $\mathcal{T}$, aljabar-$\Bbbk$ komutatif $B$, dan
	funktor serat $\omega:\mathcal{T}\to\cated{Mod}B$, tuliskan
	$\Aut^{\otimes}(\omega)$ untuk grup automorfisme monoidal $\omega$. Maka,
	untuk setiap aljabar-$B$ komutatif $B'$, terdapat isomorfisme kanonik
	\[ \Aut^{\otimes}\left( \omega \dotimes{B} B' \right) \simeq \Hom_{B\dcate{CAlg}}\left(\coEnd_B(\omega), B' \right), \]
	dengan perkalian grup pada ruas kanan tercermin oleh koperkalian $\coEnd_B(\omega)$.
\end{corollary}

Lebih lanjut, dapat ditunjukkan bahwa unsur identitas dan operasi invers dalam grup automorfisme masing-masing berasal dari morfisme kounit dan antipoda $\coEnd_B(\omega)$. Pemeriksaannya sederhana dan formal, sehingga diserahkan sebagai latihan bab ini. Dengan meninjau pembenaman Yoneda dari $B\dcate{CAlg}$, diperoleh kesimpulan berikut:
\[\begin{tikzcd}
	\text{mengetahui funktor }\Aut^{\otimes}\left( \omega \dotimes{B} (\cdot)\right): B\dcate{CAlg} \to \cate{Grp} \arrow[Leftrightarrow, d] \\
	\text{mengetahui aljabar Hopf-$B$ }\coEnd_B(\omega).
\end{tikzcd}\]
Dengan demikian, begitu $\Aut^{\otimes}$ dipandang secara funktorial, seluruh kisah panjang mengenai koaljabar endomorfisme akhirnya kembali kepada grup automorfisme-$\otimes$ dari funktor serat.

Dalam situasi Proposisi \ref{prop:Isom-Tannakian}, Deligne membuktikan dalam
\cite[1.11--1.13]{Del90} bahwa $\coHom_B(\omega_1,\omega_2)$ datar setia sebagai
modul-$B$; tentu saja ini juga menyiratkan bahwa
$\coHom_B(\omega_1,\omega_2)$ taknol. Kita tidak akan memberikan pembuktian
terperinci, tetapi perlu menyebutkan arti pentingnya dalam geometri aljabar.
Hasil tersebut menjamin adanya aljabar-$B$ datar setia $B'$ dan isomorfisme
$\omega_1\dotimes{B}B'\simeq\omega_2\dotimes{B}B'$; misalnya, ambil
$B'=\coHom_B(\omega_1,\omega_2)$. Jadi, funktor serat unik hingga perubahan
gelanggang datar setia, sedangkan sifat-sifat di atas gelanggang asal $B$ pada
prinsipnya dapat direkonstruksi melalui data penurunan dari
\S\ref{sec:descent-modules}. Inilah manfaat utama kedataran setia.
